\documentclass[11pt,letterpaper]{article}
\usepackage{fontspec}
\usepackage[margin=1in]{geometry}
\usepackage{amsmath,amssymb,amsthm,mathtools}
\usepackage{unicode-math}
\setmainfont{Latin Modern Roman}
\setmonofont{Latin Modern Mono}
\setmathfont{Latin Modern Math}
\usepackage{microtype,booktabs,array,enumitem}
\usepackage[hyphens]{url}
\usepackage[hidelinks]{hyperref}
\usepackage{fancyhdr}
\pagestyle{fancy}\fancyhf{}
\fancyhead[L]{\small Prefix stabilizer maxima}\fancyhead[R]{\small Manuscript version 2}
\fancyfoot[C]{\thepage}\renewcommand{\headrulewidth}{0pt}
\setlength{\headheight}{14pt}\setlength{\emergencystretch}{2em}
\newtheorem{theorem}{Theorem}
\newtheorem{lemma}[theorem]{Lemma}
\newtheorem{proposition}[theorem]{Proposition}
\theoremstyle{definition}\newtheorem{definition}{Definition}
\newcommand{\Stab}{\operatorname{Stab}}
\newcommand{\Max}{\operatorname{Max}}
\newcommand{\ind}{\mathbf{1}}
\newcommand{\D}{\mathcal D}
\newcommand{\C}{\mathcal C}
\newcommand{\hash}[1]{{\small\ttfamily\nolinkurl{#1}}}
\setlist[enumerate]{itemsep=3pt,topsep=5pt}
\title{\bfseries Prefix stabilizer maxima for a prime power group times a cyclic prime power}
\author{Research manuscript\quad Version 2}
\date{11 October 2026}
\begin{document}\maketitle\thispagestyle{plain}
\begin{abstract}
Let $H$ be an arbitrary finite $q$-group of order $q^a$, and let $p\ne q$ be prime. For a labeled ordering of $H\times C_{p^b}$, sum the orders of the left setwise stabilizers of all nonempty prefixes, subtracting one at each prefix. For all $a,b\ge0$, the maximum is
\[
q^ap^b\left[a\left(1-\frac1q\right)+b\left(1-\frac1p\right)\right].
\]
The maximizers are exactly the recursive consecutive right-coset orders through maximal subgroups. In every nontrivial case, the proper-prefix stabilizers recover the top subgroup uniquely, yielding an exact labeled counting recurrence. A numerical checkpoint hierarchy retaining every enclosing higher-level clip proves the upper bound at arbitrary cyclic depth. Whole-gap densification and an exact quotient identity reduce it to a symbolic $q$-adic edge inequality. If $a,b\ge1$, an order whose proper-prefix stabilizers generate the whole group loses at least $p^{b-1}(p-1)$. This sufficient gap is not claimed sharp. The proof covers both prime-power axes and the singleton; finite controls and the separate stronger historical $D_8\times C_p$ result are distinguished from the unrestricted argument.
\end{abstract}

\paragraph{Version and review status.}
Version 2 extends the earlier eleven-page $H\times C_p$ working note to every cyclic exponent $b\ge0$. It is a new manuscript and leaves the original source, PDF, and sealed package unchanged. The all-depth frozen mathematical theorem has received independent project acceptance; this is not journal acceptance. The exact acceptance record and source identities are bound in Appendix~\ref{app:provenance}. Document fidelity is a separate review gate. The bounded historical comparison performed for $b=1$ is retained with its original scope and is not a new priority audit of the all-depth theorem.

\section{Statement and conventions}
Let $q,p$ be distinct primes, let $|H|=n=q^a$ with $H$ a finite $q$-group, and let $P_b=C_{p^b}$, where $a,b\ge0$ and $P_0=1$. Put $G=H\times P_b$ and $N=np^b$. For an ordering $\mathbf{x}=(x_1,\ldots,x_N)$ of the labeled elements, define
\begin{align}
T_j&=\{x_1,\ldots,x_j\},& S_j&=\{g\in G:gT_j=T_j\},\label{eq:prefix}\\
Z_G(\mathbf{x})&=\sum_{j=1}^{N}(|S_j|-1),& B(\mathbf{x})&=\langle S_1,\ldots,S_{N-1}\rangle.\label{eq:score}
\end{align}
The action is by left multiplication. The full prefix is included in $Z_G$ and excluded from $B$. Write $M(K)$ for the maximum score on a finite group $K$ and $A(K)$ for the number of labeled maximizing orders. No quotient by translations or automorphisms is taken.

\begin{definition}
A recursive maximal-subgroup right-coset ordering of a nontrivial group chooses a maximal proper subgroup $J$, places its distinct right cosets consecutively, and recurses independently inside each coset. The singleton has its unique ordering. The identification $Jx\to J$, $hx\mapsto h$, defines the local recursion. Changing representatives merely changes these coordinates by right translation, which preserves left stabilizers.
\end{definition}

\begin{theorem}\label{thm:main}
For $G=H\times P_b$ as above, set
\begin{equation}
V_{a,b}=np^b\left[a\left(1-\frac1q\right)+b\left(1-\frac1p\right)\right].\label{eq:value}
\end{equation}
Then:
\begin{enumerate}
\item $M(G)=V_{a,b}$.
\item An ordering maximizes $Z_G$ if and only if it is a recursive maximal-subgroup right-coset ordering.
\item For nontrivial $G$, the top subgroup of every maximizing ordering is unique and equals $B(\mathbf{x})$.
\item With $A(1)=1$, for $b\ge1$,
\begin{equation}
A(H\times P_b)=p!\,A(H\times P_{b-1})^p
 +q!\sum_{M\in\Max(H)}A(M\times P_b)^q.\label{eq:count}
\end{equation}
The sum is empty when $H=1$. For $b=0$ and $H\ne1$,
\begin{equation}
A(H)=q!\sum_{M\in\Max(H)}A(M)^q.\label{eq:qcount}
\end{equation}
\item If $a,b\ge1$ and $B(\mathbf{x})=G$, then
\begin{equation}
Z_G(\mathbf{x})\le V_{a,b}-p^{b-1}(p-1).\label{eq:escape}
\end{equation}
\end{enumerate}
The singleton has score $0$, one ordering, and no top subgroup. No strict escape assertion is made for it or for either prime-power axis.
\end{theorem}

The result allows $p=2<q$ and does not require $H$ to be abelian or to have pairwise permutable subgroups. Its second factor must be cyclic. It claims no extension to a noncyclic $p$-group, to three distinct prime factors, or to all finite nilpotent groups. Neither sharpness of~\eqref{eq:escape} nor worldwide historical priority is asserted. The earlier stronger escape estimate for $D_8\times C_p$ concerns $b=1$ and is retained separately in Section~\ref{sec:d8}.

\section{Splitting and all checkpoint levels}\label{sec:checkpoints}
Every subgroup $S\le H\times P_b$ splits as $L\times Q$. For any $(h,z)\in S$, coprime-order powers isolate $(h,1)$ and $(1,z)$; doing this for every element proves the decomposition. Since $P_b$ is cyclic, $Q$ is one of its unique subgroups $P_r$ of order $p^r$, $0\le r\le b$.

A prefix is a level-$s$ checkpoint when it is $P_s$-invariant, equivalently a union of complete $P_s$-fibres. Let $C_s$ be the set of its physical indices, including the empty and full indices. Use every checkpoint at each level, not a selected subsequence. Thus
\begin{equation}
C_b\subseteq C_{b-1}\subseteq\cdots\subseteq C_1\subseteq C_0=\{0,\ldots,N\},
\quad C_s\subseteq p^s\mathbb Z,\quad 0,N\in C_s.\label{eq:hierarchy}
\end{equation}
For $j>0$, put
\[
r(j)=\max\{s:j\in C_s\},\qquad \beta_j=|S_j\cap H|.
\]
Then $|S_j|=p^{r(j)}\beta_j$, and $\beta_j$ is a power of $q$. For
\[
\D=\{q,q^2,\ldots,q^a\},\qquad c_d=d-d/q,
\]
telescoping gives
\begin{equation}
\beta-1=\sum_{\substack{d\in\D\\d\mid\beta}}c_d,
\qquad \gcd(t,n)=1+\sum_{\substack{d\in\D\\d\mid t}}c_d.\label{eq:weights}
\end{equation}
The empty set $\D$ is allowed when $n=1$.

\section{The all enclosing clipping invariant}\label{sec:clips}
First derive the necessary clips for an actual ordering. Suppose $r=r(j)<b$. For each $s>r$, let $a_s<j<z_s$ be the consecutive $C_s$ indices enclosing $j$. The inequalities are strict since $j\notin C_s$. If $d\in\D$ divides $\beta_j$, then
\begin{equation}
d\mid j/p^r,\qquad
p^sd\left\lceil\frac{a_s}{p^sd}\right\rceil<j
 <p^sd\left\lfloor\frac{z_s}{p^sd}\right\rfloor
\quad\text{for every }s>r.\label{eq:clips}
\end{equation}
To see this, quotient the $P_r$-invariant prefix by $P_r$. Its size is $j/p^r$ in $H\times(P_b/P_r)$. Among the fibres of $P_s/P_r$, its sets of full and nonempty fibres are invariant under the $H$ part of $S_j$. This action on $H\times(P_b/P_s)$ is free, so those cardinalities are multiples of $d$. The enclosing endpoints have fibre ranks $a_s/p^s$ and $z_s/p^s$. If the full and nonempty fibre counts are $f,e$, nesting gives
\[
d\left\lceil\frac{a_s}{p^sd}\right\rceil\le f,
\qquad e\le d\left\lfloor\frac{z_s}{p^sd}\right\rfloor,
\qquad p^{s-r}f\le j/p^r\le p^{s-r}e.
\]
Equality at either outer bound forces every occupied fibre to be full and hence makes the prefix $P_s$-invariant, a contradiction. Freeness on the $P_r$ quotient gives the divisibility in~\eqref{eq:clips}. For a top checkpoint $j=p^bk$, the corresponding subset of $H$ has size $k$, so $\beta_j\mid\gcd(k,n)$.

Now allow \emph{any numerical hierarchy} $\C=(C_0,\ldots,C_b)$ satisfying~\eqref{eq:hierarchy}, whether realizable or not. For $j>0$, with $r=r(j)$, define
\begin{equation}
h_{\C}(j)=
\begin{cases}
\gcd(j/p^b,n),&r=b,\\
1+\displaystyle\sum_{d\in\D}c_d I_{\C}(d,j),&r<b,
\end{cases}\label{eq:cap}
\end{equation}
where $I_{\C}(d,j)=1$ exactly when $d\mid j/p^r$ and all of~\eqref{eq:clips} hold, for every $s>r$. Otherwise it is zero. Put
\begin{equation}
U_b(\C)=\sum_{j=1}^{N}\left[p^{r(j)}h_{\C}(j)-1\right].\label{eq:invariant}
\end{equation}
The cap need not be realizable or a subgroup order. For the actual checkpoint hierarchy,~\eqref{eq:weights}--\eqref{eq:clips} imply $\beta_j\le h_{\C}(j)$. Retain the entire top-endpoint deficit
\begin{equation}
D_{\rm top}=\sum_{\substack{k>0\\p^bk\in C_b}}
 p^b\left[\gcd(k,n)-\beta_{p^bk}\right].\label{eq:deficit}
\end{equation}
Its summands are nonnegative and its final summand is zero. Thus the strengthened upper bound is
\begin{equation}
Z_G(\mathbf{x})\le U_b(\C)-D_{\rm top}.\label{eq:strongcap}
\end{equation}
The induction will preserve the numerical hierarchy, every enclosing clip, and this unchanged deficit from the actual top checkpoints.

\section{Whole gap densification}\label{sec:dense}
Assume $b\ge1$. Let $pu<pv$ be consecutive $C_1$ indices and write $R=v-u\ge1$. There is no positive-level checkpoint in their open physical interval. Therefore for every $s\ge2$, its enclosing $C_s$ endpoints $a_s,z_s$ are constant throughout this interval. For $d\in\D$ put
\begin{align}
L_d&=\max\left(\{u\}\cup\left\{p^{s-1}d\left\lceil\frac{a_s}{p^sd}\right\rceil:2\le s\le b\right\}\right),\label{eq:L}\\
U_d&=\min\left(\{v\}\cup\left\{p^{s-1}d\left\lfloor\frac{z_s}{p^sd}\right\rfloor:2\le s\le b\right\}\right).\label{eq:U}
\end{align}
When $b=1$, these are just $u,v$. All higher-level clipping endpoints in these expressions are multiples of $d$. If $L_d\ge U_d$, define $N_d=R_d=0$. Otherwise set
\begin{equation}
N_d=\max\left(0,p\left[\left\lfloor\frac{U_d}{d}\right\rfloor-
\left\lceil\frac{L_d}{d}\right\rceil\right]-1\right),
\qquad R_d=\#\{t\in\mathbb Z:L_d<t<U_d,\ d\mid t\}.\label{eq:NR}
\end{equation}
The old open interval contributes exactly $\sum_d c_dN_d$ to $U_b$. Its indices all have level zero. Eligibility for $d$ requires divisibility by $d$ and the strict level-one clip
\[
pd\lceil u/d\rceil<j<pd\lfloor v/d\rfloor.
\]
Intersecting with all higher clips gives
\[
pd\lceil L_d/d\rceil<j<pd\lfloor U_d/d\rfloor,
\]
whose eligible multiples of $d$ are counted by $N_d$.

Densify $C_1$ by inserting \emph{all} missing inner ranks $pt$, $u<t<v$, at once, while leaving every $C_s$, $s\ge2$, unchanged. The inserted indices have level one and their eligibility for $d$ is exactly $d\mid t$ and $L_d<t<U_d$. Every remaining level-zero index has no $q$-weight in its cap: it lies between adjacent new $C_1$ ranks, too close to enclose the required distinct $d$-multiples for $d\ge q\ge2$. Its cap $h_{\C}$ is therefore $1$ and its score contribution is zero. The new open-interval contribution is
\[
(p-1)(R-1)+p\sum_d c_dR_d.
\]
The new minus old difference is
\begin{equation}
\Delta=(p-1)(R-1)+\sum_d c_d(pR_d-N_d).\label{eq:delta}
\end{equation}
Here is a complete arithmetic check of its sign. For $L_d\ge U_d$, the summand is zero. Otherwise there are three cases:
\begin{enumerate}
\item Both endpoints are divisible by $d$. If $w=(U_d-L_d)/d\ge1$, then $R_d=w-1$, $N_d=pw-1$, and $pR_d-N_d=-(p-1)$.
\item Exactly one endpoint is divisible by $d$. Then $N_d=\max(0,pR_d-1)$; the difference is $0$ for $R_d=0$, and $1$ otherwise.
\item Neither endpoint is divisible by $d$. The difference is $0$ for $R_d=0$, $p$ for $R_d=1$, and $p+1$ for $R_d\ge2$.
\end{enumerate}
Only the first case is negative, and it requires $d\le U_d-L_d\le R$. Put
\[
d_* = \max\{q^i:0\le i\le a,\ q^i\le R\}.
\]
In particular $d_*\le\min(R,n)$, even for gaps with $R>n$. The telescoping identity gives
\[
\sum_{\substack{d\in\D\\d\le R}}c_d=d_*-1\le R-1,
\]
and hence
\begin{equation}
\Delta\ge(p-1)\bigl[R-1-(d_*-1)\bigr]=(p-1)(R-d_*)\ge0.\label{eq:deltasign}
\end{equation}
All pre-existing checkpoint caps stay unchanged, since a positive-level cap only uses levels strictly above its own. Repeating over the gaps, let $\C^{\rm dense}$ have $C_1$ equal to all multiples of $p$. Then
\begin{equation}
U_b(\C)\le U_b(\C^{\rm dense}).\label{eq:densification}
\end{equation}
This compares whole gaps of numerical bounds. It assumes neither single-checkpoint insertion monotonicity nor realizability of the densified hierarchy by a group ordering.

\section{Exact quotient identity at arbitrary depth}\label{sec:quotient}
In $\C^{\rm dense}$, every nonmultiple of $p$ has zero score contribution. At each remaining index $j=pt$, decrement its level by one and divide higher checkpoint indices by $p$. The resulting hierarchy for total length $np^{b-1}$ is
\begin{equation}
C'_0=\{0,\ldots,np^{b-1}\},\qquad
C'_s=\{j/p:j\in C_{s+1}\}\quad(1\le s\le b-1).\label{eq:quotient}
\end{equation}
The clipping, divisibility, and top-$H$ caps correspond exactly. The retained $q$-power data remain $\D=\{q,\ldots,q^a\}$. At each positive quotient index,
\[
p^rh-1=p\bigl[p^{r-1}h-1\bigr]+(p-1),
\]
so there is an exact identity
\begin{equation}
U_b(\C^{\rm dense})=pU_{b-1}(\C')+(p-1)np^{b-1}.\label{eq:quotientidentity}
\end{equation}
For depth zero, the definition is $U_0=\sum_{k=1}^{n}[\gcd(k,n)-1]$. Equation~\eqref{eq:quotientidentity} is an identity of the numerical invariant; it does not invoke a group theorem on an artificial quotient order.

The intermediate length $np^{b-1}$ may have both prime divisors. It is never substituted for $n$ in $\gcd(t,n)$ or in $\D$. The $H$ rank base remains the prime power $n=q^a$; only the cyclic $p$-depth decreases.

Apply~\eqref{eq:densification} and~\eqref{eq:quotientidentity} $b-1$ times. Every additive term, after preceding scalings, is $(p-1)np^{b-1}$. Therefore
\begin{equation}
U_b(\C)\le p^{b-1}U_1(\C^{\rm top})+(b-1)(p-1)np^{b-1},\label{eq:collapse}
\end{equation}
where $C^{\rm top}_1=\{j/p^{b-1}:j\in C_b\}$. The actual top checkpoints never change. The actual deficit $D_{\rm top}$ in~\eqref{eq:strongcap} remains subtractable after every replacement; no endpoint deficit is discarded and no earlier stabilizer is assumed to persist at later prefixes.

The same elimination restricts to a consecutive top interval $p^bk<p^bl$. Its additive term is
\[
(b-1)(p-1)p^{b-1}(l-k).
\]
This interval form will give the exact outer edge loss.

\section{The outer q adic edge and upper bound}\label{sec:edge}
Put $h(k)=\gcd(k,n)$, $s(k)=\sum_{t=1}^{k}h(t)$, and $s(0)=0$. For $0\le k<l\le n$ define
\begin{align}
w_d(k,l)&=\max\left(0,\left\lfloor\frac ld\right\rfloor-\left\lceil\frac kd\right\rceil\right),\label{eq:w}\\
E_p(k,l)&=p\left[h(l)+\sum_{d\in\D}c_dw_d(k,l)\right]
-\left[1+\sum_{\substack{d\in\D\\w_d(k,l)>0}}c_d\right],\label{eq:edge}\\
F_p(k)&=ps(k)-k.\label{eq:potential}
\end{align}
At one level, the strict clips count $\max(0,pw_d-1)$ eligible open-interval indices for $d$, and the complete right endpoint contributes $ph(l)-1$. Thus $U_1$ is exactly the sum of $E_p$ over consecutive top ranks. This includes composite prime powers $n=q^a$ and the empty $\D$ case.

For completeness, the one-level calculation is reproduced. Put $m=l-k$, $g=\gcd(k,l,n)$, and let $X$ consist of those $d\in\D$ for which neither endpoint is divisible by $d$ but an interior $d$-multiple exists. Then
\begin{align}
F_p(l)-F_p(k)-E_p(k,l)&=p\alpha+\gamma,\label{eq:loss}\\
\alpha&=m-g+\sum_{d\in X}c_d,\label{eq:alpha}\\
\gamma&=1-m+\sum_{\substack{d\in\D\\w_d>0}}c_d.\label{eq:gamma}
\end{align}
Indeed, if $N_d$ counts strict interior $d$-multiples, expanding $h(t)$ gives
\[
\alpha=m-1+\sum_d c_d(N_d-w_d).
\]
The difference $N_d-w_d$ is $-1$ for two divisible endpoints; it is $+1$ for two indivisible endpoints with an interior multiple; it is zero otherwise. The negative weights sum to $g-1$, proving~\eqref{eq:alpha}.

Since $g\mid m$, if $m>g$ then $m\ge2g$, and
\begin{align}
p\alpha+\gamma&\ge p(m-g)+1-m\notag\\
&=(p-1)(m-2g)+(p-2)g+1\ge p-1.\label{eq:positive}
\end{align}
If $m=g$, the interval is aligned $q$-adic: its endpoints are consecutive multiples of the $q$-power $m$. For $d\le m$, both endpoints are divisible and $w_d=m/d$; for $d>m$, there is no interior $d$-multiple. Consequently $X$ is empty, $\alpha=\gamma=0$, and
\begin{equation}
F_p(l)-F_p(k)-E_p(k,l)\ge0,
\quad\text{with equality exactly when }l-k=\gcd(k,l,n).\label{eq:zero}
\end{equation}
Every nonzero loss is at least $p-1$, including $p=2$.

For $b\ge1$ define
\begin{align}
E_b(k,l)&=p^{b-1}E_p(k,l)+(b-1)p^{b-1}(p-1)(l-k),\label{eq:collapsededge}\\
F_b(k)&=p^{b-1}F_p(k)+(b-1)p^{b-1}(p-1)k\notag\\
&=p^bs(k)+\bigl[bp^{b-1}(p-1)-p^b\bigr]k.\label{eq:collapsedpotential}
\end{align}
Equation~\eqref{eq:collapse} and its interval form give
\begin{equation}
Z_G\le\sum_{\text{consecutive top ranks }k,l}E_b(k,l)-D_{\rm top}.\label{eq:outerbound}
\end{equation}
The exact loss is
\begin{equation}
F_b(l)-F_b(k)-E_b(k,l)=p^{b-1}\bigl[F_p(l)-F_p(k)-E_p(k,l)\bigr].\label{eq:scaledloss}
\end{equation}
Thus it has the same zero criterion, and every nonzero loss is at least $p^{b-1}(p-1)$. No unproved propagation of losses through new layers is needed.

Counting the $q$-adic valuations of $1,\ldots,n$ gives
\[
s(n)=n+\sum_{i=1}^{a}(q^i-q^{i-1})\frac{n}{q^i}
=n+an\left(1-\frac1q\right).
\]
The outer potential telescopes from $0$ to $n$, and hence
\begin{equation}
Z_G\le F_b(n)=V_{a,b}.\label{eq:upper}
\end{equation}
This establishes the unrestricted bound for every $b\ge1$. The $b=0$ group case is proved independently in Section~\ref{sec:bases}.

\section{Strict top escape}\label{sec:escape}
Assume $a,b\ge1$ and $B(\mathbf{x})=G$. There must be a proper top checkpoint. Otherwise every proper-prefix stabilizer has $p$-part contained in the unique $P_{b-1}$, and the subgroup they generate lies in $H\times P_{b-1}<G$. This step depends on the cyclic factor: several smaller projections cannot together generate $P_b$.

Set $r_0=n/q$. A top-rank path consisting entirely of zero-loss edges either is the single edge $0\to n$ or visits $r_0$. Indeed, every edge in a nontrivial zero-loss path is an aligned $q$-adic interval of length $d<n$, so $d\le n/q$ and $d\mid r_0$. An interval between consecutive multiples of $d$ cannot contain the multiple $r_0$ strictly inside it.

The proper top checkpoint makes the path nontrivial. If it omits $r_0$, it has a nonzero outer edge, and~\eqref{eq:scaledloss} gives~\eqref{eq:escape}.

If $r_0$ is present and $\beta_{p^br_0}=r_0$, its subset of $H$ is one right coset of its stabilizer $K$ of order $r_0$. This $K$ is normal of index $q$ (Lemma~\ref{lem:normal} below). The whole checkpoint is therefore one right coset of $J=K\times P_b$, normal of prime index $q$ in $G$. Any element outside $J$ acts in a single cycle on the $q$ quotient cosets. Before this checkpoint, a nonempty prefix lies in that one $J$-coset and its translate lies in a disjoint coset. After this checkpoint, if such an element stabilized a prefix, its iterates would force all translates of the full coset into that prefix, hence all of $G$. Thus every proper stabilizer lies in $J$, contradicting $B(\mathbf{x})=G$.

It follows that a present $r_0$ must have $\beta_{p^br_0}<r_0$. Its contribution to $D_{\rm top}$ is at least $p^b>p^{b-1}(p-1)$, so~\eqref{eq:outerbound} again gives~\eqref{eq:escape}. When $r_0=1$, this deficient case is impossible because $\beta\ge1$. No stabilizer persistence between different prefixes is used.

\section{Prime power bases and local accounting}\label{sec:bases}
We prove the complete prime-power dependency, including both axes $b=0$ and $a=0$.

\begin{lemma}\label{lem:normal}
Every maximal proper subgroup of a nontrivial finite $t$-group is normal of index $t$.
\end{lemma}
\begin{proof}
Every proper subgroup $L<Q$ of a finite $t$-group is strictly contained in its normalizer. The class equation gives a nontrivial center $Z(Q)$. If $Z(Q)\not\subseteq L$, a central element outside $L$ belongs to $N_Q(L)\setminus L$. Otherwise induction applies to $L/Z(Q)<Q/Z(Q)$ and its strict normalizer lifts to one for $L$. If $L$ is maximal, $N_Q(L)=Q$. The quotient $Q/L$ is a nontrivial $t$-group with no nontrivial proper subgroup, so it is cyclic of prime order $t$.
\end{proof}

\begin{lemma}[Normal prime index block accounting]\label{lem:block}
Suppose $J\lhd K$ has prime index $t$. If an ordering consists of consecutive right $J$-cosets, then every proper-prefix stabilizer is contained in $J$ and equals the stabilizer of its active local prefix in $J$. Its score is
\begin{equation}
Z_K=\sum_C Z_J(\mathbf{x}|_C)+|K|-|J|.\label{eq:block}
\end{equation}
\end{lemma}
\begin{proof}
An element outside $J$ acts transitively on the $t$ quotient cosets. Inside a block, it cannot fix the unique partially occupied coset. At a proper block boundary, it cannot stabilize the nonempty proper set of full quotient cosets. Thus every proper stabilizer lies in $J$. Elements of $J$ preserve every full and empty coset, so precisely the active local stabilizer remains. Every local contribution appears once, except that the final contribution $|J|-1$ is replaced by $|K|-1$.
\end{proof}

The following identity also applies to interleaved local orders. Suppose $S_j\le J$ for every $j<|K|$. At insertion $j$ into $C_j=Jx_j$, let $L_j^{\rm loc}=\Stab_J(T_j\cap C_j)$. Then $S_j\le L_j^{\rm loc}$ and
\begin{equation}
Z_K=\sum_C Z_J(\mathbf{x}|_C)+|K|-|J|
-\sum_{j<|K|}\bigl(|L_j^{\rm loc}|-|S_j|\bigr).\label{eq:local}
\end{equation}
Each local prefix is counted at its active insertion. The last local and global stabilizers are $J$ and $K$, respectively; every displayed deficit is nonnegative.

If equality forces all these deficits to vanish, the first completion of a $J$-coset makes the entire prefix $J$-invariant. No other coset has yet completed; invariance forces all others to be empty. That coset is the first consecutive block. At each subsequent completion before the last, after the previously completed blocks, the same argument forces exactly the next consecutive block. The final coset fills the remainder. This is the first-completion argument used below.

Now let $|Q|=t^e$. A prefix stabilizer acts freely and its order divides both $j$ and $t^e$. Therefore
\begin{equation}
Z_Q\le\sum_{j=1}^{t^e}\bigl[\gcd(j,t^e)-1\bigr]
=t^e e\left(1-\frac1t\right).\label{eq:basevalue}
\end{equation}
For $Q=1$, the value is $0$ and there is one ordering. For nontrivial $Q$, choose a maximal subgroup $J$, normal of index $t$ by Lemma~\ref{lem:normal}. Optimal local orders exist by induction on $e$, and Lemma~\ref{lem:block} gives attainment of~\eqref{eq:basevalue}.

Conversely, equality in~\eqref{eq:basevalue} saturates every individual prefix cap. At $j=|Q|/t$, the prefix is one right coset of its stabilizer $J$ of that order. The prime-index coset trap from Section~\ref{sec:escape} forces every proper stabilizer into $J$. The local identity~\eqref{eq:local} and inductive value force every local order optimal and every deficit zero. The first-completion argument then supplies consecutive coset blocks. Induction proves exactly the recursive characterization in both directions. At the first full block the stabilizer is $J$, and all proper stabilizers lie in $J$, so $B(\mathbf{x})=J$. Hence the top is unique, and
\[
A(Q)=t!\sum_{J\in\Max(Q)}A(J)^t\quad(Q\ne1),\qquad A(1)=1.
\]
Apply this to $Q=H$ when $b=0$ and to $Q=P_b$ when $a=0$. It proves both axes and their singleton intersection without a general nilpotent theorem.

\section{Both maximal subgroup directions and equality}\label{sec:equality}
Induct on $a+b$, with both axes supplied by Section~\ref{sec:bases}. For $a,b\ge1$ the numerical upper bound is already~\eqref{eq:upper}. Choose $J=H\times P_{b-1}$, normal of index $p$. Its local theorem has smaller $a+b$. Consecutive $J$-cosets with optimal local orders have score
\begin{equation}
pV_{a,b-1}+np^b-np^{b-1}=V_{a,b},\label{eq:attain}
\end{equation}
so the bound is attained.

For any maximizing order, strict escape implies $B(\mathbf{x})<G$. Choose a maximal proper subgroup $J$ containing it. Splitting and the cyclic subgroup chain show that all maximal subgroups are exactly
\begin{equation}
H\times P_{b-1}\quad\text{of index }p,
\qquad M\times P_b\quad(M\in\Max(H))\quad\text{of index }q.\label{eq:maxtypes}
\end{equation}
A split subgroup with both factors proper can enlarge either and is not maximal. With only the $p$-factor proper, maximality requires $P_{b-1}$; with only the $H$-factor proper, it requires a maximal $M$. Every subgroup in~\eqref{eq:maxtypes} is normal of the indicated prime index. The local theorem is the $(a,b-1)$ case, the $(a-1,b)$ case, or a prime-power base. These are all the proper-subgroup inputs required.

All proper-prefix stabilizers lie in $J$, so~\eqref{eq:local} applies. With $t=[G:J]$, the local value satisfies
\begin{equation}
tM(J)+|G|-|J|=V_{a,b}.\label{eq:bothvalues}
\end{equation}
For index $p$ this is~\eqref{eq:attain}; for index $q$ it is
\[
qV_{a-1,b}+np^b-(n/q)p^b=V_{a,b}.
\]
Equality of the global score forces every local order optimal and every local deficit zero. The first-completion argument proves consecutive $J$-coset blocks; local optimality and induction give the entire recursive hierarchy.

Conversely, any recursive ordering uses one of the top subgroups in~\eqref{eq:maxtypes}. Its local orders are optimal by induction and its value is $V_{a,b}$ by Lemma~\ref{lem:block} and~\eqref{eq:bothvalues}. Every proper stabilizer lies in $J$, while at the first full block the stabilizer is exactly $J$. Thus $B(\mathbf{x})=J$, which recovers the top uniquely. This proves equality in both directions without needing to classify equality at every numerical densification step.

For fixed maximal $J$ of index $t$, the distinct right cosets have $t!$ orders, and each independently has $A(J)$ labeled optimal internal orders. Fixing right-translation coordinates gives $t!A(J)^t$ possibilities; changing representatives introduces no additional factor. Unique top recovery makes different $J$ classes disjoint. Summing over~\eqref{eq:maxtypes} proves~\eqref{eq:count}; the prime-power proof supplies~\eqref{eq:qcount}. For $H=1$, the sum is empty, and $A(P_b)=p!A(P_{b-1})^p$. For the singleton, $A(1)=1$ and no top subgroup is asserted. This completes Theorem~\ref{thm:main}.

\section{Why nearest level clipping does not suffice}\label{sec:diagnostic}
The use of every enclosing clip is essential to this proof method. For $n=8$, $q=2$, $p=3$, $b=3$, consider the numerical hierarchy
\[
C_1=C_2=C_3=\{0,27,54,\ldots,216\}.
\]
In the first outer interval, through index $27$, a weakened cap testing only $s=r(j)+1$ has value $55$. Densifying $C_1$ lowers that weakened cap to $48$, contradicting the needed monotonic direction. Explicitly, the old open cap is $11\cdot1+5\cdot2+2\cdot4=29$, with endpoint cap $26$. After densification, eight proper level-one nodes give baseline $16$ and $q$-weight increment $6$, for $16+6+26=48$.

Under the actual all-enclosing invariant, the level-three clip suppresses all old open $q$-weights. The old cap is $26$, the first densification gives $16+26=42$, and subsequent elimination reaches $54$. The correct comparison is preserved. This is a counterexample to the weakened proof method, not to the group theorem; no group realizability of this numerical hierarchy is assumed.

\section{The stronger historical dihedral result}\label{sec:d8}
In this section only, let $G=D_8\times P$ and $P=C_p$, so $b=1$. The earlier independently accepted theorem for $D_8\times C_p$, with $p$ odd, proves
\begin{equation}
B(\mathbf{x})=G\quad\Longrightarrow\quad
Z_G(\mathbf{x})\le18p-7=(20p-8)-(2p-1).\label{eq:d8escape}
\end{equation}
Here $V_{3,1}=20p-8$. This specialized loss $2p-1$ is stronger than the general loss $p-1$; the present theorem does not replace or weaken it. Its separate proof uses the complete symbolic 36-edge table on ranks $0,\ldots,8$ and the index-two checkpoint trap. The unchanged theorem and acceptance review are included among the historical source records. Equation~\eqref{eq:d8escape} is reported as that prior result, not derived from~\eqref{eq:escape}.

For its explicit labeled count, put $f=p!$ and
\[
T=f\,2^p+2f^2,\qquad U=f\,8^p+2T^2,\qquad V_4=f\,24^p+6T^2.
\]
These are respectively the counts for $C_2\times C_p$, $C_4\times C_p$, and $C_2^2\times C_p$. Since $D_8$ has one cyclic and two Klein-four maximal subgroups,
\[
A(D_8)=2\cdot8^2+4\cdot24^2=2432,
\]
and the historical recurrence gives
\begin{equation}
A(D_8\times C_p)=p!\,2432^p+2U^2+4V_4^2.\label{eq:d8count}
\end{equation}
At $p=3$, the value is $52$ and the count is $203{,}047{,}403{,}520$.

There is also a uniform escaping example. Write
$D_8=\langle r,s:r^4=s^2=1,\ srs=r^{-1}\rangle$. The order
\[
e,s,r,rs,r^2,r^2s,r^3,r^3s
\]
has proper-prefix stabilizers generating $D_8$ and score $10$. Expanding each element to its full $P$-fibre gives $B(\mathbf{x})=D_8\times P$ and
\[
Z_G=p\cdot10+8(p-1)=18p-8.
\]
This example is one below the specialized upper bound. Neither that observation nor the general theorem asserts the sharp escaping maximum. It also shows why unconditional accumulated-coset containment cannot be used outside its proved hypotheses.

\section{Finite controls separate from the proof}\label{sec:controls}
Sections~\ref{sec:clips}--\ref{sec:equality} give a symbolic proof for unrestricted $a,b$ and all distinct primes $p,q$. Finite enumeration is not a premise. The frozen author and independent review packets contain separately written corroborating controls; they were inspected and source-bound for this presentation, not executed during typesetting.

The author packet reports $336$ arbitrary numerical hierarchies through $b=4$, $4{,}325$ whole-gap comparisons, $336$ exact quotient and outer-reduction checks, and $36{,}459$ one-level outer edges with scaled checks through $b=5$. Its $33$ direct group-order checks reach order $216$, including $D_8\times C_{27}$ and the order-$27$ Heisenberg group times $C_8$. Constructed escaping orders cover both the missing-$r_0$ and deficient-$r_0$ branches.

The independent acceptance packet reports $1{,}015$ hierarchies, $1{,}011$ quotient identities, $15{,}582$ whole-gap comparisons including $1{,}650$ gaps longer than $n$, and $155{,}520$ outer edges through depth $b=8$, with no negative-loss failure. It checks $170$ direct orders across $D_8\times C_{27}$, $C_2^3\times C_{27}$, the order-$27$ Heisenberg group times $C_8$, $C_9\times C_{16}$, and $C_{32}$. Three orders have $B=G$, including both escape branches for $D_8\times C_{27}$. The weakened nearest-only diagnostic and the correct all-enclosing comparison are independently reproduced.

An independent dynamic program over every subset uses
\[
D(S)=|\Stab(S)|-1+\max_{x\in S}D(S\setminus\{x\}),\qquad D(\varnothing)=0,
\]
with counts summed over maximizing last elements. It does not assume the optimizer hierarchy. Table~\ref{tab:controls} records the four cyclic controls from the all-depth review. The counts agree with separately calculated recursive-family counts. These checks establish finite corroborations only; they do not imply any unrestricted quantifier or cover noncyclic second Sylow factors.

\begin{table}[htbp]\centering
\caption{Independent exhaustive cyclic controls in the all-depth review}\label{tab:controls}
\begin{tabular}{lrrr}\toprule
Group & Subset states & Maximum & Maximizing labeled orders\\\midrule
$C_6$ & $64$ & $7$ & $120$\\
$C_{12}$ & $4{,}096$ & $20$ & $31{,}872$\\
$C_{18}$ & $262{,}144$ & $33$ & $13{,}727{,}232$\\
$C_{24}$ & $16{,}777{,}216$ & $52$ & $2{,}044{,}231{,}680$\\\bottomrule
\end{tabular}\end{table}

\section{Related work and the contribution boundary}\label{sec:related}
The mathematical increment in this version is the upper bound and necessity of the full optimizer hierarchy for $H\times C_{p^b}$ at arbitrary cyclic depth. Unique top recovery and the labeled recurrence are consequences of that structure. Independent correctness review and historical comparison answer different questions. The comparison summarized here was performed for the earlier $b=1$ result; it has not been repeated as a priority audit of the larger family.

For a common labeling, a regular-action permutation $g$ has an invariant proper initial segment exactly at a connectivity-set cut in the terminology of Stanley~\cite{stanley}. Hence $Z_G$ is the sum of connectivity-set sizes over all group elements, including the identity's $|G|-1$ cuts. The present optimization concerns simultaneous relabelings of the entire fixed regular group, preserving multiplication relations. The inspected connectivity statements in the bounded comparison do not supply the earlier $H\times C_p$ maximum or its equality cases.

Fixed-chain transversal signatures are classical; Gonz\'alez Vasco, R\"otteler, and Steinwandt~\cite{gvs} give such constructions and minimum-length signatures for solvable groups. Their length objective is the sum of factor-block sizes. Coset-dependent constructions also predate this work: Baumeister and de Wiljes~\cite{bdw}, Algorithm~5.2 and Theorem~5.4, combine local signatures indexed by transversal choices under stated compatibility conditions. Independent branching or coset-dependent choice is therefore not claimed as an original idea here. Those inspected statements concern factorization blocks, rather than the score over every nested element prefix, and were not found to supply the earlier prefix upper bound or all-optimizer necessity theorem.

Likewise the sum over top choices, factorial ordering of children, and product over independent subproblems are standard counting mechanisms. Maximal-chain enumeration and rooted-tree recursions are relevant antecedents~\cite{st,mbkk}. The objective-specific input to~\eqref{eq:count} is identifying all optimizers and proving disjointness across top subgroups; no new counting principle is claimed.

The bounded $b=1$ comparison located no exact earlier equivalent. That result cannot be promoted to an all-depth priority conclusion. The comparison with Erd\H{o}s, Faigle, and Kern~\cite{efk} remains incomplete: Theorem~2.1 and its main incidence calculation were recovered, but full definitions and Lemma~2.1, most of Section~4, the proof of Proposition~5.1, and the complete Section~6 ending remain partly unread. The checked direct substitutions fail the required hypotheses; this is not a whole-paper exclusion. The preserved historical report records the exact search and access boundaries. No worldwide-first claim, sharp escape result, or general nilpotent theorem follows from this manuscript.

\appendix
\section{Version 2 change index}\label{app:changes}
The preceding eleven-page working note concerned $H\times C_p$. This manuscript is a separate version and makes the following source-bound changes.
\begin{enumerate}
\item The statement replaces $b=1$ by every $b\ge0$, replaces the value and labeled recurrence accordingly, and restricts the sufficient escape gap to $a,b\ge1$. Both axes and the singleton are explicit.
\item Sections~\ref{sec:checkpoints}--\ref{sec:quotient} introduce every checkpoint level, the all-enclosing invariant, the full actual top deficit, the whole-gap comparison, and the exact affine quotient identity. The $q$-power base $n$ is fixed throughout the reduction.
\item Sections~\ref{sec:edge}--\ref{sec:escape} retain the symbolic one-level $q$-adic edge calculation and scale its loss exactly by $p^{b-1}$. They use the unique cyclic subgroup $P_{b-1}$ and the prime-index coset trap to prove strict escape.
\item Sections~\ref{sec:bases}--\ref{sec:equality} supply the prime-power dependencies and induction on $a+b$ in both maximal-subgroup directions, proving attainment, all equality cases, unique top recovery, and labeled counting.
\item Section~\ref{sec:diagnostic} records the depth-three obstruction to nearest-only clipping. Section~\ref{sec:controls} distinguishes new all-depth finite controls from the symbolic proof.
\item Section~\ref{sec:d8} preserves the stronger historical $D_8\times C_p$ result only at $b=1$. Section~\ref{sec:related} preserves the bounded historical comparison and makes no new all-depth priority assertion.
\end{enumerate}
The accompanying source-change index maps these sections to the frozen theorem and prior version. The original files and accepted mathematical sources are copied unchanged; the manuscript re-typesets them without changing their status or original version identities.

\section{Exact source identities}\label{app:provenance}
The sealed package binds the following sources by exact byte count and SHA-256. Original mathematical programs were not executed during manuscript preparation. The new package's delivery receipt records its Library round-trip separately to avoid a self-referential archive hash. Document fidelity review is recorded separately from the mathematical acceptance below.

\paragraph{All-depth author archive.}
Library version 0 has identifier\begin{center}\hash{libfile_1dae9a88f5108191a048f86017d447d6}\end{center}The archive is $42{,}673$ bytes, SHA-256:
\begin{center}\hash{f41867be27ab0f77376b6745adff35575aa3af5044ce4f29d201f60179858d3e}\end{center}
Its unchanged theorem is $21{,}879$ bytes, SHA-256:
\begin{center}\hash{e13674de3c970bc3a18417d9f6e01630b388a0632571637e97674c526f56afd9}\end{center}
The archive retains its original candidate label; acceptance is established by the separate exact-byte review.

\paragraph{Independent all-depth acceptance.}
Library version 0 has identifier\begin{center}\hash{libfile_672d090b20a4819198c511cd125dd489}\end{center}The archive is $80{,}379$ bytes, SHA-256:
\begin{center}\hash{efddb76d4d402aebedf75e1c095ed7c22776ab90142966c0a3a64049cb94f62d}\end{center}
The decision is \emph{ACCEPT as written}, with no mathematical correction. The review is $19{,}175$ bytes, SHA-256:
\begin{center}\hash{05443fa967c2eee7694c64fd1dc8719bfe3fbe9c9c3bad21b897cfc56a262066}\end{center}
The exact accepted statement is $2{,}611$ bytes, SHA-256:
\begin{center}\hash{eceb8796e9f1018c44d596dc8024cce6d6abf34211800259a63d9537d3e037a9}\end{center}
It accepts precisely $H\times C_{p^b}$ for distinct primes $p,q$ and $a,b\ge0$, with the guards in Theorem~\ref{thm:main}. It does not accept historical priority or a broader group class.

\paragraph{Original version preserved.}
The unchanged eleven-page PDF has Library identifier\begin{center}\hash{libfile_78c45eb171c08191b2f53963f2aa4845}\end{center}It is $104{,}885$ bytes, SHA-256:
\begin{center}\hash{24448fdfa5652d6c6c8a182cab38c035cd39d96bba37e1b2abf6f5e65c884b19}\end{center}
The original editable LaTeX has Library identifier\begin{center}\hash{libfile_ea8ce3f6559481919dfddff2ec859e21}\end{center}It is $32{,}673$ bytes, SHA-256:
\begin{center}\hash{cb9cd8f829ce0efc1bec90eed3f5ef896224238515be25fc204024055aacb9d0}\end{center}
The original package has Library identifier\begin{center}\hash{libfile_d2e3566fd3e08191b6c589e14cd43813}\end{center}It is $406{,}953$ bytes, SHA-256:
\begin{center}\hash{245121a1567f6dc4bf35b2b29b775dadae108e968e3bebb1a3e67f0a42cac36d}\end{center}

\paragraph{Historical comparison preserved at its bounded scope.}
The Library identifier is\begin{center}\hash{libfile_07a6103f05bc8191bb3dde2b717f4fc7}\end{center}The archive is $12{,}425$ bytes, SHA-256:
\begin{center}\hash{f6bffb4b177ba1cb067ce09765af64ff8c853f3fdebd59f3653354c9abd37ada}\end{center}
This is the $H\times C_p$ historical comparison, not a mathematical proof premise or an all-depth priority audit. The source bindings also hash the historical $D_8$ theorem and review, the complete accepted statement, and the independently reported finite controls. Third-party full texts are not included in the new package.

\begin{thebibliography}{9}
\bibitem{stanley} R. P. Stanley, \emph{The Descent Set and Connectivity Set of a Permutation}. Primary manuscript dated 16 August 2005, equation~(1), Proposition~1.1, and Theorems~1.2 and~1.4. \url{https://math.mit.edu/~rstan/papers/decset.pdf}.
\bibitem{gvs} M. I. Gonz\'alez Vasco, M. R\"otteler, and R. Steinwandt, \emph{On Minimal Length Factorizations of Finite Groups}, Experimental Mathematics 12 (2003). Definition~2.1 and Proposition~3.1. \url{https://www.microsoft.com/en-us/research/wp-content/uploads/2016/11/On-Minimal-Length-Factorizations-of-Finite-Groups.pdf}.
\bibitem{bdw} B. Baumeister and J.-H. de Wiljes, \emph{Aperiodic logarithmic signatures}, arXiv:1104.4503 (2011), Algorithm~5.2 and Theorem~5.4. \url{https://arxiv.org/abs/1104.4503}.
\bibitem{st} M. \c{S}tef\u{a}nescu and M. T\u{a}rn\u{a}uceanu, \emph{Counting maximal chains of subgroups of finite nilpotent groups}, Carpathian Journal of Mathematics 25 (2009), inspected pages 119--124, Theorem~2.1 and Section~3. \url{https://www.carpathian.cunbm.utcluj.ro/wp-content/uploads/2009-vol-25-1/carpathian_2009_25_1_119_127.pdf}.
\bibitem{mbkk} F. A. Matsen IV, S. C. Billey, A. Kas, and M. Konvalinka, \emph{Tanglegrams A Reduction Tool for Mathematical Phylogenetics}, primary manuscript, Section~2.2, Theorem~5 and Example~6. \url{https://matsen.fredhutch.org/papers/Matsen2018-bg.pdf}.
\bibitem{efk} P. Erd\H{o}s, U. Faigle, and W. Kern, \emph{A Group-Theoretic Setting for Some Intersecting Sperner Families}, Combinatorics, Probability and Computing 1 (1992), 323--334. Partially inspected; the remaining access boundary is stated in Section~\ref{sec:related}. \url{https://doi.org/10.1017/S0963548300000377}.
\end{thebibliography}
\end{document}
